Il seguente Capitolo tratterà l'implementazione di un Processo Leggero nel Linguaggio di Programmazione Java.

\section{Creazione di un Processo Leggero}
Ogni Processo Leggero in Java viene rappresentato da un Oggetto di Tipo \texttt{Thread}\cite{Thread_JavaDocs} o di un suo Sotto-Tipo: questa semantica permette di creare Classi che possano essere eseguite su \textit{Thread} separati dal principale, permettendo quindi un migliore parallelismo nell'esecuzione delle loro operazioni. \par
Java mette a disposizione del Programmatore due diverse possibilità per la stesura di un Programma che utilizzi più Processi Leggeri: l'estensione della Classe \texttt{Thread} e l'implementazione dell'Interfaccia \texttt{Runnable}. \par
Prima di descrivere entrambe le opportunità è bene fornire informazioni sull'ereditarietà della Classe \texttt{Thread}: essa infatti estende la Classe \texttt{Object} e implementa l'interfaccia \texttt{Runnable} così come raffigurato nel Diagramma UML di seguito.
\begin{figure}[H]
    \centering
    \includesvg[width=.25\textwidth]{images/assets/JavaThreadInheritance.svg}
    \caption[Eredirarietà della Classe \texttt{Thread}]{Eredirarietà della Classe \texttt{Thread} in Java. UML Generato tramite l'ausilio dell'editor PlantText\cite{PlantText}.}
    \label{fig:ThreadInheritance}
\end{figure}
Come affermato in precedenza, l'obiettivo ultimo di un Programma che utilizzi più Processi Leggeri è quello di avere un oggetto di tipo \texttt{Thread} (o di un suo Sotto-Tipo) per ognuno di essi, così da sfruttarne le Primitive. \par
Si illustrano ora le due diverse possibilità di stesura di un Programma Multi-Processo.
\subsection{Estensione della Classe \texttt{Thread}}
Un primo approccio alla Programmazione \textit{Multi-Threaded} in Java può essere quello di Estendere la Classe \texttt{Thread}\cite{Thread_JavaDocs} con la propria. \par
L'unico requisito è quello di completare l'implementazione del Metodo \texttt{run()} della Super-Classe esplicitando le operazioni da eseguire una volta iniziata l'esecuzione del nuovo \textit{Thread}. \par
Si riporta a seguire un esempio didattico per l’utilizzo dove il Thread principale andrà a creare un numero prestabilito di Processi Leggeri, i quali dovranno stampare sullo Standard Output un messaggio con l'identificativo loro assegnato in fase di inizializzazione per poi terminare l’esecuzione.
Nel \textit{Main} verrà creato un oggetto di tipo \texttt{ClassExtendThread} per ogni Processo Leggero, per poi utilizzarne la primitiva \texttt{start()} ereditata da \texttt{Thread}.
\begin{lstlisting}[language=Java, numbers=left, stepnumber=1]
public class ClassExtendThread extends Thread {
    
    // Definisco l'Identificativo che verra' assegnato dal Main
    private int threadNumber;

    public ClassExtendThread(int tNum) {
        // Inizializzo l'Identificativo
        this.threadNumber = tNum;
    }

    // Implemento il metodo run() sovrascrivendo l'implementazione di Thread
    @Override
    public void run() {
        // Stampo l'Identificativo
        System.out.println("Sono il Thread numero " + this.threadNumber + ".");
    }
}   
\end{lstlisting}
\captionof{lstlisting}{Esempio di Classe che estende la Classe Thread per essere eseguita in un Processo Leggero separato.}
\subsection{Implementazione dell'Interfaccia \texttt{Runnable}}
Alternativa alla modalità precedente è quella di implementare l'Interfaccia \texttt{Runnable}\cite{Runnable_JavaDocs}. \par
Tale cambiamento non modificherà in maniera radicale il codice, ma necessiterà di un differente ragionamento: ora infatti andrà creato un oggetto di tipo \textit{Thread} a partire da un oggetto della Classe implementante \texttt{Runnable}, per poi operare direttamente sull'oggetto di tipo \textit{Thread}. \par
Si riporta a seguire una reimplementazione dell'esempio precedente seguendo questa nuova logica.
\begin{lstlisting}[language=Java, numbers=left, stepnumber=1]
public class ClassImplementRunnable implements Runnable {
    ...
}
\end{lstlisting}
\captionof{lstlisting}{Modifica effettuata al precedente esempio per adattare il codice al nuovo approccio.}
Com'è possibile notare dal Codice sopra le modifiche all'interno della Classe che implementi l'Interfaccia \texttt{Runnable} saranno minime: i requisiti sono infatti esattamente gli stessi. Il vero cambiamento si avrà nella classe principale che dovrà gestire come segue la creazione dei nuovi Processi Leggeri.
\begin{lstlisting}[language=Java, numbers=left, stepnumber=1]
...
for(int i=0; i<NUM_OF_THREADS; i++) {
    // Creo un nuovo oggetto di tipo ClassImplementRunnable
    ClassImplementRunnable cir = new ClassImplementRunnable(i);
    // Creo un Thread a partire dall'oggetto della mia Classe
    Thread thread = new Thread(cir);
    // Eseguo il Thread
    thread.start();
}
...
\end{lstlisting}
\captionof{lstlisting}{Classe Main per l'utilizzo della Classe precedentemente illustrata.}
\subsection{Considerazioni Pratiche ed Esecuzione}
Le due implementazioni, seppur molto simili visivamente, risultano differenti da un punto di vista logico per le motivazioni indicate poc'anzi. Per ragioni pratiche legate al fatto che in Java non risulta possibile estendere più di una Classe alla volta ma è possibile implementare più Interfacce contemporaneamente, si preferisce generalmente procedere con il secondo approccio. \par
L'esecuzione dei due Programmi così come sono stati stesi risulterà nel medesimo risultato: una stampa in ordine randomico da parte di ciascuno dei Processi Leggeri invocati.
\begin{lstlisting}[language=Bash]
$ java MainImplementRunnable
Sono il Main e sto per creare i Thread.
Sono il Main ed ho terminato l'esecuzione.
Sono il Thread numero 2.
Sono il Thread numero 3.
Sono il Thread numero 0.
Sono il Thread numero 4.
Sono il Thread numero 1.
\end{lstlisting}
\captionof{lstlisting}{Esempio di esecuzione dei Programmi sopra riportati.}
Si noti come l'esecuzione del \textit{Thread} principale termini prima di quella degli altri Processi Leggeri da lui generati ma che questa circostanza non limiti in alcun modo questi ultimi, che continueranno la loro esecuzione fino alla rispettiva terminazione.

\section{Primitive per la Gestione dell'Esecuzione}
La classe \texttt{Thread} mette a disposizione Metodi per la gestione dell'Esecuzione dei singoli Processi Leggeri\cite{SeminarioJavaDispense}. Ciò risulta utile nel caso in cui si voglia influenzare il loro ordine di esecuzione o ancora la loro possibilità di eseguire. \par
\begin{itemize}
    \item \textbf{Il metodo \texttt{start()}} \par
    Il metodo \texttt{start()} permette di inzializzare l'esecuzione del Processo Leggero invocando il metodo \texttt{run()} della Classe.
    \item \textbf{Il metodo \texttt{stop()}} \par
    Il metodo \texttt{stop()} permette di interrompere invece l'esecuzione del Processo Leggero sollevando un'Eccezione di tipo \texttt{ThreadDeath}.
    \item \textbf{Il metodo \texttt{suspend()}} \par
    Il metodo \texttt{suspend()} permette di sospendere momentaneamente l'esecuzione di un \textit{Thread} fin tanto che su di lui non venga chiamato il metodo \texttt{resume()}, portandolo dallo stato \textit{Running} allo stato \textit{Blocked}.
    \item \textbf{Il metodo \texttt{resume()}} \par
    Il metodo \texttt{resume()} permette di riprendere l'esecuzione di un \textit{Thread} bloccato dal metodo \texttt{suspend()}, portandolo dallo stato \textit{Blocked} allo stato \textit{Ready}.
    \item \textbf{Il metodo \texttt{sleep(int millis)}} \par
    Il metodo \texttt{sleep()} permette di portare il Processo Leggero nello stato \textit{Blocked} per un determinato numero di millisecondi senza utilizzare cicli di Clock della CPU. Tale metodo solleva un'Eccezione di tipo \texttt{InterruptedException} nel caso in cui a tale Processo Leggero venisse chiesto di interrompersi mentre è già in stato \textit{Blocked}, perciò il metodo va incapsulato in un costrutto di tipo \textit{Try-Catch}.
    \item \textbf{Il metodo \texttt{destroy()}} \par
    Il metodo \texttt{destroy()} permette di eliminare il \textit{Thread} senza effettuare alcuna pulizia, mantenendo quindi i vari \textit{Lock} da esso utilizzati e facilitando l'insorgere di una situazione di Blocco Critico in quanto nessun altro Processo potrà più accedere a tali risorse.
    \item \textbf{Il metodo \texttt{yield()}} \par
    Il metodo \texttt{yield()} permette al \textit{Thread} di lasciare volontariamente la CPU ad un altro processo, implicando un successivo \textit{Re-Scheduling} dei Processi da eseguire.
    \item \textbf{Il metodo \texttt{join()}} \par
    Il metodo \texttt{join()} consente al Processo Padre di attendere la terminazione del Processo Figlio prima di continuare con la propria esecuzione, garantendo un semplice meccanismo di sincronizzazione tra i due.
\end{itemize}
\begin{figure}[H]
    \centering

    \includesvg[width=\textwidth]{images/assets/ThreadJoinFlow.svg}

    \caption[Flusso di Esecuzione alterato dal Metodo ``\texttt{join()}''.]{Flusso di Esecuzione di un Processo Padre ed un Processo Figlio nel caso di utilizzo del Metodo ``\texttt{join()}''. Grafo Generato tramite l'ausilio dell'editor PlantText\cite{PlantText}.}
    \label{fig.ThreadJoinExecutionFlow}
\end{figure}
\subsection{Considerazioni}
Come discusso in precedenza, l'arbitraria interruzione dell'esecuzione di un Processo potrebbe far incorrere in situazioni indesiderate: le risorse da lui utilizzate potrebbero rimanere infatti in uno stato inconsistente o indefinitamente bloccate se tale Processo aveva acquisito un \textit{Lock} per l'accesso esclusivo a tale risorsa facilitando l'insorgere di situazioni di Blocco Critico, infine la possibilità di bloccare per un tempo arbitrario l'esecuzione di un \textit{Thread} potrebbe portare alla \textit{Starvation} di questo. \par
Per questi motivi i Metodi \texttt{stop()}, \texttt{suspend()}, \texttt{resume()} e \texttt{destroy()} sono stati deprecati a partire da Java 1.2 nel 1998 e se ne scoraggia fortemente l'utilizzo. \par
Alle primitive sopracitate si preferisce usufruire della capacità di segnalazione Inter-Processo comunicando la possibilità di un \textit{Thread} di sospendere o terminare la propria esecuzione, in maniera tale da permettergli di eseguire la propria Sezione Critica e lasciare il Sistema in uno stato consistente. Un esempio di Metodo che permette al Processo Leggero di lasciare volontariamente la CPU è \texttt{yield()}. La natura delle moderne versioni della \textit{Java Virtual Machine} implica poi l'utilizzo dello \textit{Scheduler} del Sistema Operativo, il quale dovrà scegliere il prossimo Processo cui lasciare il controllo della CPU. Notando che tutti i moderni Algoritmi di \textit{Scheduling} supportano l'utilizzo di un valore di Priorità per la scelta del Processo dalla Coda dei Processi Pronti, anche Java mette a disposizione un valore di Priorità per i propri Processi Leggeri: esso varia tra il valore minimo ``\texttt{MIN\textunderscore PRIORITY}'' (0) ed il valore massimo ``\texttt{MAX\textunderscore PRIORITY}'' (10) e viene impostato in maniera automatica - al netto di ulteriori istruzioni - su ``\texttt{NORM\textunderscore PRIORITY}'' (5), risultando visualizzabile e personalizzabile rispettivamente tramite le Primitive \texttt{getPriority()} e \texttt{setPriority()}. Tra le altre primitive utilizzate per una migliore sincronizzazione tra i processi risulta indubbiamente presente il Metodo \texttt{join()} che permete al Processo Padre di attendere la terminazione del Processo Figlio prima di continuare con la propria Esecuzione.
\subsection{Esempio di Controllo dell'Esecuzione tramite \texttt{join()} e \texttt{yield()}}
Si riporta a seguire il Codice di un semplice Programma generante due \textit{Thread} che incrementano un contatore correlato di un esempio di esecuzione.
\begin{lstlisting}[language=Java, numbers=left, stepnumber=1]
public class ThreadExecutionControl {
    public static void main(String[] args) {
        // Creo un'istanza della Classe di Supporto con il Lock
        ThreadExecutionControlCounter tecc = new ThreadExecutionControlCounter();
        // Creo due Thread
        Thread t1 = new ThreadExecutionControlThread(1, tecc);
        Thread t2 = new ThreadExecutionControlThread(2, tecc);
        // Avvio i due Thread
        t1.start();
        t2.start();
        // Attendo i due Thread e stampo il risultato
        try {
            t1.join();
            t2.join();
            System.out.println("Sono il Main, il valore finale del contatore e' " + tecc.getCounter());
        } catch(InterruptedException ie) {
            System.out.println("I Thread sono stati interrotti inaspettatamente");
            ie.printStackTrace();
        }
    }
}

class ThreadExecutionControlCounter {
    // Inizializzo una variabile contatore
    private int counter = 0;
    // Incremento il contatore
    protected void increaseAndPrintCounter(int threadNum) {
        counter++;
        try {
            Thread.sleep(100);
        } catch(InterruptedException ie) {
            System.out.println("Il Thread " + threadNum + " e' stato interrotto inaspettatamente");
            ie.printStackTrace();
        }
        System.out.println("Il contatore vale " + this.counter + " ed e' stato aggiornato dal Thread " + threadNum);
    }
    // Getter del contatore
    protected int getCounter() {
        return this.counter;
    }
}

class ThreadExecutionControlThread extends Thread {
    // Identificatore del Thread
    int threadNum;
    // Oggetto della Classe di supporto
    ThreadExecutionControlCounter tecc;
    // Costruttore
    public ThreadExecutionControlThread(int threadNum, ThreadExecutionControlCounter tecc) {
        this.threadNum = threadNum;
        this.tecc = tecc;
        System.out.println("Thread " + this.threadNum + " avviato");
    }
    // Implementazione del Metodo run
    @Override
    public void run() {
        for(int i=0; i<5; i++) {
            // Incremento il contatore
            this.tecc.increaseAndPrintCounter(threadNum);
            // Rilascio volontariamente la CPU
            Thread.yield();
        }
    }
}
\end{lstlisting}
\captionof{lstlisting}{Codice Java del Programma sopra descritto.}
\begin{lstlisting}[language=Bash]
$ java ThreadExecutionControl
Thread numero 1 avviato
Thread numero 2 avviato
Il contatore vale 2 ed e' stato aggiornato dal Thread 2
Il contatore vale 1 ed e' stato aggiornato dal Thread 1
Il contatore vale 3 ed e' stato aggiornato dal Thread 2
Il contatore vale 3 ed e' stato aggiornato dal Thread 1
Il contatore vale 5 ed e' stato aggiornato dal Thread 1
Il contatore vale 4 ed e' stato aggiornato dal Thread 2
Il contatore vale 6 ed e' stato aggiornato dal Thread 2
Il contatore vale 6 ed e' stato aggiornato dal Thread 1
Il contatore vale 7 ed e' stato aggiornato dal Thread 2
Il contatore vale 7 ed e' stato aggiornato dal Thread 1
Sono il Main, il valore finale del contatore e' 7
\end{lstlisting}
\captionof{lstlisting}{Esecuzione del Programma.}
Si noti che nell'esecuzione il risultato finale della variabile contatore non è quello desiderato, e che durante l'esecuzione vengono ripetuti alcuni valori del contatore. Queste inconsistenze sono assolutamente previste date le considerazioni sulla concorrenza nell'accesso alle risorse condivise effettuate in precedenza, e per questo vengono introdotti meccanismi di Sincronizzazione tra Processi.